\capituloPregunta{¿Cómo decidir cuando hay covariables o varias respuestas?}

\section{Motivación}

No todos los experimentos se explican con una sola respuesta ni con grupos perfectamente comparables. A veces una variable adicional modifica la respuesta. Otras veces se miden varias respuestas que juntas describen el desempeño del sistema.

\section{ANCOVA}

ANCOVA permite comparar tratamientos ajustando por covariables. La pregunta cambia de:

\begin{quote}
¿Los grupos son diferentes?
\end{quote}

a:

\begin{quote}
¿Los grupos siguen siendo diferentes si consideramos una variable que también influye en la respuesta?
\end{quote}

\section{MANOVA}

MANOVA se utiliza cuando existen varias variables respuesta. El material del curso incluye ejemplos con crecimiento de plantas y con productos donde varias respuestas deben analizarse de manera conjunta.

\begin{idea}
Si varias respuestas describen el mismo sistema, analizarlas por separado puede fragmentar la decisión.
\end{idea}

\section{Actividad}

Diseñe un experimento con una variable respuesta principal y una covariable. Después proponga una versión multirrespuesta del mismo problema.

\begin{cuidado}
No use ANCOVA para ajustar variables que son consecuencia del tratamiento. No use MANOVA solo para acumular muchas variables sin una pregunta común.
\end{cuidado}

\section{Conexión}

El cierre natural del curso consiste en unir pregunta, diseño, datos, análisis, discusión y comunicación.
